Write
\[
 v_d=\mathcal V_d^*(P),\qquad
 \widehat v_d=\widehat{\mathcal V}_{d,m}^*,\qquad
 v_*=\min_{d\in\mathcal D}v_d.
\]
Choose any \(d^*\in D^*(P)\).  By the definition of \(\varepsilon_m\) and the empirical minimizing property of \(\widehat d_m\),
\[
 v_{\widehat d_m}
 \leq \widehat v_{\widehat d_m}+\varepsilon_m
 \leq \widehat v_{d^*}+\varepsilon_m
 \leq v_{d^*}+2\varepsilon_m
 =v_*+2\varepsilon_m.                                             \tag{1}
\]
Since \(v_*\) is the population minimum, (1) proves the deterministic regret inequality.

If \(\widehat d_m\notin D^*(P)\), the definition of \(\gamma(P)\) gives
\[
 \gamma(P)\leq v_{\widehat d_m}-v_*\leq2\varepsilon_m.            \tag{2}
\]
This proves the asserted event inclusion.  When \(D^*(P)=\mathcal D\), the event on the left is empty and, under the convention \(\gamma(P)=+\infty\), so is the event on the right.  Otherwise, finiteness of \(\mathcal D\) implies \(\gamma(P)>0\).

It remains to establish the stochastic claims.  The pilot construction and the budget-profile theorem give
\[
 \widehat v_d
 =F_d(\widehat V_{E,d},\widehat V_{O,d}),\qquad
 v_d=F_d(V_{E,d},V_{O,d}),
\]
where
\[
 F_d(x,y)=\bigl(\sqrt{c_{E,d}x}+\sqrt{c_{O,d}y}\bigr)^2.           \tag{3}
\]
The nondegeneracy assumption, positive costs, and finiteness of \(\mathcal D\) place all population arguments of (3) in a common compact subset of \((0,\infty)^2\).  The uniform pilot error is \(O_P(r_{V,m})=o_P(1)\), so with probability tending to one the estimated arguments also lie in a slightly larger compact subset of \((0,\infty)^2\).  The gradients of the finitely many maps \(F_d\) are uniformly bounded there.  The mean-value theorem consequently yields
\[
 \varepsilon_m
 =\max_{d\in\mathcal D}|\widehat v_d-v_d|
 =O_P(r_{V,m})=o_P(1).                                             \tag{4}
\]
Together with (2) and the strict positivity of \(\gamma(P)\) when a nonoptimal design exists, (4) proves
\[
 \Pr\{\widehat d_m\notin D^*(P)\}
 \leq\Pr\{2\varepsilon_m\geq\gamma(P)\}\longrightarrow0.       \tag{5}
\]

For the allocation claim, define
\[
 A_d(x,y)=
 \frac{\sqrt{c_{E,d}x}}
 {\sqrt{c_{E,d}x}+\sqrt{c_{O,d}y}}.
\]
The inequality
\[
 |(x\vee\underline v_m)-V|
 \leq |x-V|+\underline v_m
\]
and the same compact-set mean-value argument give
\[
 \max_{d\in\mathcal D}
 |\widehat\alpha_{d,m}-\alpha_d^*(P)|
 =O_P(r_{V,m}+\underline v_m).                                    \tag{6}
\]
Finally, the second-order allocation-loss lemma gives the stronger uniform bound
\[
\begin{aligned}
 \max_{d\in\mathcal D}\left|B\left\{
 \frac{V_{E,d}}{\lfloor\widehat\alpha_{d,m}B/c_{E,d}\rfloor}
 +\frac{V_{O,d}}{\lfloor(1-\widehat\alpha_{d,m})B/c_{O,d}\rfloor}
 \right\}-v_d\right|
 =O_P\left((r_{V,m}+\underline v_m)^2+B^{-1}\right).
\end{aligned}                                                       \tag{7}
\]
Because \(r_{V,m}+\underline v_m=o(1)\), the right side of (7) is also
\(O_P(r_{V,m}+\underline v_m+B^{-1})=o_P(1)\).  This proves every stated rate.  The argument used only the fact that the deterministic tie rule returns an empirical minimizer; (5) correctly concerns membership in the entire population minimizer set.